\documentclass[a4paper]{article}
\usepackage{luatexja}
\usepackage{amsmath}
\usepackage{amssymb}
\usepackage{geometry}
\usepackage{hyperref}

\geometry{margin=25mm}

\title{ねじれ星による重力レンズ：\\
捕獲された光線の帰還}

\author{https://github.com/hypernumbernet}
\date{\today}

\begin{document}

\maketitle

\begin{abstract}
ねじれ星はシュヴァルツシルトの外縁を保ち、したがって光子球を保つ。捕獲は保たない。許容錐の上界は、計量係数が符号を変える前に、赤方偏移を正の床へ凍結する。三軸の読みでは \(1<r_\star/r_s<1.005\)、一軸の飽和では \(r_\star/r_s=1/(1-\mathrm{e}^{-\pi})\) であり、どちらも \(r_{\mathrm{ph}}=\frac32 r_s\) の内側にある。\(b>b_c=3\sqrt{3}\,GM/c^2\) の光線は、その球の外で転向し、節に出会わない。\(b<b_c\) の光線は、ブラックホールが捕獲する光線である。古典チャートの上では転向点がなく、\(r_s\) へ降りるキリング時間は発散し、\(r<r_s\) では動径方向が時間的になっている。凍結チャートの上では、同じ光線が \(r=b\sqrt{A_{\min}}\) で、凍結の内側において転向し、戻る。光子球から凍結面までの往復は \(11.44\,GM/c^3\) または \(20.96\,GM/c^3\) であり、光子の半周 \(3\pi\sqrt{3}\,GM/c^3\) の両側に落ちる。凍結の内側で費やす時間は、それより長い。その最小値は、すでに半周を超える。一軸の飽和では \(20\,GM/c^3\) より大きく、三軸の上界では \(280\,GM/c^3\) より大きい。中心まで達する動径光線は、\(80\) より大きく、\(400\) より大きい。等スケールの節は、この有限な内部の上にある。節を符号つきの薄い殻と読むと、引力の殻は \(R>b_m/\sqrt{2}\) で内部に動径臨界曲線を加え、内部の接線曲線は \(b_m/\sqrt{2}<R<b_m\) のときに限って、動径曲線の外側に加える。斥力の殻は、自身の円盤の上にどちらの曲線も加えない。殻の質量は未定であるから、倍率は予測しない。太陽の光、銀河の環、M87* と Sgr A* の放射環は \(r_\star\) の外にとどまる。遅いエコー列は、この帰還ではない。
\end{abstract}

\section{はじめに}
\label{sec:introduction}

二重時空理論は共有連続体を拒む。各質量粒子は通常のミンコフスキー枠と双対枠を運び、重力はその対のねじれ不整合である~\cite{dst-main}。球対称質量の外縁では、不整合は動径ブーストであり、そのブーストはシュヴァルツシルト・コフレームを再現する。十分な密度の内部では、同じ不整合が半径とともに符号を変える。親理論は、その結果として得られる対象をねじれ星と呼ぶ。事象地平もなく中心特異点もない、層状の凝縮体である。層は単一の干渉因子の零点であり、\(\pi\) の各区間に節が一つある。シュヴァルツシルト・チャートへ貼り足した内部解ではない。

光が問うのは、質量粒子の落下とは別の問いである。書かれた作用は光子に消える剛性を与えるので、光子のラピディティを加速しない。外縁で光を偏向するのは、コフレームのヌル幾何である。本稿が問うのは、その幾何が、ブラックホールなら捕獲する光線をどう扱うかである。

シュヴァルツシルト・レンズは、強場では臨界インパクトパラメータを一つ、すなわち光子球をもち、弱場では臨界曲線を一つ、すなわち全質量のアインシュタイン環をもつ。どちらも凍結の外で定まるので、どちらも生き延びる。光子球を通過する光線は生き延びない。古典チャートの上では、転向する場所がない。地平へ降りるキリング時間は発散し、地平の内側では動径座標が時間的になっている。凍結チャートは、それらの光線に近点と有限な帰還を与える。ねじれ星の内部を特徴づけるのは、その帰還であって、アインシュタイン環の第二の写しではない。

二つの読みは分けておかねばならない。外縁の偏向、準地平に対する光子球の位置、節の窓、そして臨界より下の光線の帰還は、チャートの読みである。節を符号つきの薄い殻と同一視することは、障壁の弱偏向における模型である。各障壁へ質量を割り当てるものではなく、強場の内部を貫くヌル積分でもない。

第\ref{sec:chart}節はチャートを記録する。第\ref{sec:capture}節は、\(b_c\) の内側の光線をシュヴァルツシルト幾何がどう扱うかを述べる。第\ref{sec:return}節は、凍結がそれに代えて何をするかを述べる。第\ref{sec:shells}節は節を弱いレンズとして読む。第\ref{sec:observations}節は外縁類を、太陽の遅延、銀河の環、M87* と Sgr A* の放射環と比較する。第\ref{sec:discriminants}節は、二つの理論を実際に分けうる間隔を切り出す。

\section{チャート}
\label{sec:chart}

外縁では、動径ブーストと赤方偏移は
\begin{equation}
\label{eq:exterior-A}
\varphi(r)=-\frac12\log\bigl(1-r_s/r\bigr),
\qquad
A(r)=1-\frac{r_s}{r},
\qquad
r_s=\frac{2GM}{c^2},
\end{equation}
であり、\(r>r_s\) である。コフレーム \(\bigl(\sqrt{A},\,A^{-1/2},\,r,\,r\sin\theta\bigr)\) はシュヴァルツシルト計量を誘導する。このチャートには三つの二次量が住み、それらは同一の対象ではない。有限角の不整合 \(J=\frac12\varphi^2\)、場の種 \(J_{\mathrm{field}}=\frac12(\varphi')^2\)、そしてテレパラレル・スカラー \(T\) である。弱場の主要次数では一致し、それ以外では一致しない。アインシュタイン方程式は、テトラッド変分ののち \(T\) に属する。\(J\) の恒等式ではない。

古典シュヴァルツシルトは \(r\to r_s^+\) で \(\varphi\to\infty\)、\(A\to 0\) とし、\(r<r_s\) では因子 \(A\) を負にする。許容錐は、その符号変化の前に発散を切る。三軸の上界 \(|J|=3\pi^2/8\) を動径の純粋ブーストと読むと、
\begin{equation}
\label{eq:quasihorizon}
\varphi_{\max}=\frac{\pi\sqrt{3}}{2},
\qquad
A_{\min}=\mathrm{e}^{-\pi\sqrt{3}},
\qquad
\frac{r_\star}{r_s}=\frac{1}{1-A_{\min}}
\end{equation}
へ反転し、\(0<A_{\min}<10^{-2}\)、\(1<r_\star/r_s<1.005\) である。\(r\le r_\star\) ではチャートは \(\varphi=\varphi_{\max}\) と \(A=A_{\min}\) を凍結する。凍結した因子は狭義に正であるから、ヌル面 \(A=0\) は形成されず、静的キリング場は内部を通じて時間的なままである。親理論の同じ箇所は、この同一視が三軸上界を一つのブースト軸へ読んだものであると標す。一軸の許容ブーストは \(\varphi=\pi/2\) で飽和し、そこでは \(A=\mathrm{e}^{-\pi}\)、\(r_\star/r_s=1/(1-\mathrm{e}^{-\pi})\) であり、その比もなお \(3/2\) より小さい。

層状の内部を組織するのは別のスカラーである。等スケール・アンザッツのもと、\(x=\ell/(2r)\) で書いた干渉因子は
\begin{equation}
\label{eq:gamma-s}
\gamma_s(x)=\cosh x\cos x-\sinh x\sin x
\end{equation}
であり、導関数は \(\gamma_s'(x)=-2\cosh x\,\sin x\) である。各区間 \([n\pi,(n+1)\pi]\) はちょうど一つの節を含み、\(\gamma_s\) の符号は平坦部から平坦部へと交替し、節は \(n\pi+\pi/4<x_n<n\pi+\pi/2\) にある。したがって最初の二つの面積半径は
\begin{equation}
\label{eq:radius-ratio}
\frac16<\frac{r_2}{r_1}<\frac{4}{5\pi}\approx 0.255
\end{equation}
を満たす。比は無次元の根だけに依る。長さ \(\ell\) は絶対尺度を定め、比からは相殺される。平坦部は \(1/\cosh\bigl(\ell/(2r)\bigr)\) によって遮蔽される。中心へ向かって成長するのは節の障壁である。許容錐上では \(|J|\le 3\pi^2/8\) であり、テレパラレル密度は \(-8<r^2 T<27\) にとどまるので、スパイクは有限である。錐はその窓を一意な面密度へ変換しない。

書かれた作用において、光子は剛性が消える粒子である。剛性が消えると、共通ラピディティの加速が除かれる。光子を真空へ強制することはなく、\(F=k/(\gamma_s r^2)\) を光線への力にすることもない。外縁のヌル偏向はコフレーム~\eqref{eq:exterior-A} の測地線である。障壁を通るヌル偏向は、第\ref{sec:shells}節の模型である。

\section{捕獲}
\label{sec:capture}

静的で球対称な線素
\[
ds^2=-A(r)\,c^2\,dt^2+A(r)^{-1}\,dr^2+r^2\,d\Omega^2
\]
はエネルギー \(E\) と角運動量 \(L\) を保存する。その比が漸近インパクトパラメータ \(b=L/E\) である。動径の転向点があるとすれば、それは
\begin{equation}
\label{eq:turning}
b^2=\frac{r^2}{A(r)}
\end{equation}
を満たす半径である。外縁因子について、右辺を \(u=r/r_s>1\) で書くと \(r_s^2\,u^3/(u-1)\) である。三次の恒等式
\[
4u^3-27u+27=(2u-3)^2(u+3)
\]
は、この量が少なくとも \(\frac{27}{4}r_s^2\) であり、等号が成り立つのは \(u=\frac32\) のときに限ると告げる。その最小値が光子球
\begin{equation}
\label{eq:photon-sphere}
r_{\mathrm{ph}}=\frac32 r_s=\frac{3GM}{c^2},
\qquad
b_c=\frac{r_{\mathrm{ph}}}{\sqrt{A(r_{\mathrm{ph}})}}=3\sqrt{3}\,\frac{GM}{c^2}
\end{equation}
である。\(r^2/A\) の最小は有効ポテンシャル \(A/r^2\) の最大であるから、円軌道は不安定である。無限遠から来る \(b>b_c\) の光線は、ある \(r_0>r_{\mathrm{ph}}\) に転向点をもつ。\(b<b_c\) の光線は、\(r>r_s\) の上に転向点をもたない。あらゆる古典半径が、光線の運ぶものより大きいインパクトパラメータを要求するので、光線は転向しない。

同じチャートは、内向きの行程を無限に長くする。動径ヌル光線は \(c\,dt=dr/A\) を満たし、
\begin{equation}
\label{eq:antideriv}
\frac{d}{dr}\Bigl(r+r_s\log|r-r_s|\Bigr)=\frac{1}{A(r)}
\end{equation}
が \(r>r_s\) で成り立つ。\(r\downarrow r_s\) で対数は発散するので、提案されるいかなる遅れに対しても、地平のなお外にある半径が存在し、そこでは光子球からのキリング時間がすでにその遅れを超える。光線が \(r_s\) に達するのは無限のキリング時間においてであり、光線は戻らない。

地平の内側では、障害は無限の遅れより鋭い。\(0<r<r_s\) では \(A<0\) である。\(dt^2\) の係数と \(dr^2\) の係数は符号を交換している。\(t\) は空間的であり、\(r\) は時間的である。正のインパクトパラメータの転向点は、右辺が負である \(b^2=r^2/A\) を要求するが、それは不可能である。光線は動径運動を反転する自由をもたない。古典チャートが記述するブラックホールの内部は、これである。

転向する光線の方位角の進み、すなわち転向点 \(r_0>r_{\mathrm{ph}}\) から無限遠へ出て戻る進みは
\begin{equation}
\label{eq:deflection-integral}
\Delta\varphi(b)
  =2\int_{r_0}^{\infty}
    \frac{dr}{r^2\sqrt{b^{-2}-A(r)/r^2}},
\qquad
\alpha_{\mathrm{ext}}(b)=\Delta\varphi(b)-\pi
\end{equation}
である。弱場では \(4GM/(c^2 b)\) へ還元される。同じ因子がシャピロ遅延を定める。パラメトライズド後ニュートン記法では、曲がりも遅延も \((1+\gamma)/2\) に比例し、チャート~\eqref{eq:exterior-A} は \(\gamma=1\) の場合である。

薄いレンズでは、レンズ方程式が角で測った光源位置 \(\beta\) と像位置 \(\theta\) を \(\beta=\theta-\bar\alpha(\theta)\) で結び、\(\bar\alpha=(D_{LS}/D_S)\,\alpha(D_L\theta)\) である。点質量 \(M\) のアインシュタイン角は
\begin{equation}
\label{eq:einstein}
\theta_E^2=\frac{4GM}{c^2}\,\frac{D_{LS}}{D_L D_S}
\end{equation}
であり、弱い偏向は \(\bar\alpha=\theta_E^2/\theta\) となる。レンズ写像の接線固有値が消えるのは \(\bar\alpha=\theta\) においてであり、動径固有値が消えるのは \(d\bar\alpha/d\theta=1\) においてである~\cite{NarayanBartelmann1996}。点質量は正の根 \(\theta=\theta_E\) を一つもち、動径の根をもたない。層がさらなる焦点を描くなら、ねじれ星が弱場で超えねばならないレンズはこれである。捕獲と帰還との区別は、これではない。

\section{帰還}
\label{sec:return}

\(r>r_\star\) では外縁因子~\eqref{eq:exterior-A} を、\(r\le r_\star\) では凍結値 \(A=A_{\min}\) を、つなぐ。接合は連続である。凍結区間では \(A/r^2\) は \(r\) について狭義に減少するので、凍結そのものは光子球を寄与しない。両方の上界は \(r_{\mathrm{ph}}\) の内側にある。三軸の比は \(1.005\) より小さく、一軸の比 \(1/(1-\mathrm{e}^{-\pi})\) は \(3/2\) より小さい。無限遠からの光線が越えねばならない唯一の障壁は、インパクトパラメータ~\eqref{eq:photon-sphere} にあるシュヴァルツシルトの光子球のままである。\(r_\star\) の内側の層はそれを動かさない。

したがって光線は、古典の問題と同じ二類に分かれ、そのあとで異なる。

\(b>b_c\) の光線は \(r_0>r_{\mathrm{ph}}>r_\star\) で転向する。偏向される経路の全体がシュヴァルツシルト領域にある。節には出会わない。インパクトパラメータが \(r_s\) に比べて大きいあらゆる弱場の試験、およびこの臨界曲線の角直径 \(2b_c/D_A\) を観測量とするあらゆる強場の試験は、この類に属する。

\(0<b<b_c\) の光線は、第\ref{sec:capture}節が捕獲されたと宣告した光線である。接合したチャートの上では、その宣告は成り立たない。\(r>r_s\) にはなお転向点がないので、光線は光子球を通過し、凍結へ入る。凍結の内側では \(A\) は正の定数 \(A_{\min}\) であり、~\eqref{eq:turning} は正の根を一つもつ。
\begin{equation}
\label{eq:periapsis}
r_{\mathrm{peri}}=b\sqrt{A_{\min}}.
\end{equation}
どちらの上界についても、この近点は \(r_\star\) の内側にある。\(b\) が \(b_c\) へ下から近づく極限においてさえ、そうである。光子球をわずかに越えただけの光線も、\(r_\star\) と \(r_{\mathrm{ph}}\) のあいだの空洞ではなく、凍結の内側で転向する。\(A\) は符号を変えていないので、動径方向はなお空間的であり、チャートの静的観測者はなお存在し、光線は反転する。それを吸収する地平はない。光線は \(r_{\mathrm{ph}}\) を再び横切り、無限遠へ達しうる。\(b_c\) の近くでは、落下の前に、不安定軌道の傍らを過ぎる通過が何度も巻きうる。その巻きは外縁幾何の光子環列であり、一つの臨界曲線を繰り返し訪れる。節のリムの列ではない。

行程のキリング時間は \(r_\star\) で分かれる。光子球から凍結面へ降りて戻る往復は、~\eqref{eq:antideriv} によって
\begin{equation}
\label{eq:echo-delay}
\Delta t
  =\frac{2}{c}\int_{r_\star}^{r_{\mathrm{ph}}}\frac{dr}{A(r)}
  =\frac{2}{c}\Bigl[r+r_s\log|r-r_s|\Bigr]_{r_\star}^{r_{\mathrm{ph}}}
\end{equation}
である。\(r_\star>r_s\) ならば積分は有限である。\(A_{\min}=\mathrm{e}^{-\pi\sqrt{3}}\) では \(20.96\,GM/c^3\) に等しい。一軸の飽和 \(A=\mathrm{e}^{-\pi}\) では \(11.44\,GM/c^3\) に等しい。親理論は三軸の読みを一軸へ移植した上界だと標すので、その同一視が決着するまでは、チャートが予測するのは \(\{11.44,\,20.96\}\,GM/c^3\) の中の間隔であり、単一の数ではない。

外にすでに存在する自然な時計は、不安定光子軌道の半周期である。\(r=3GM/c^2\) 上の座標角速度は \(d\varphi/dt=c^3/(3\sqrt{3}\,GM)\) であるから、
\begin{equation}
\label{eq:half-orbit}
T_{1/2}=3\pi\sqrt{3}\,\frac{GM}{c^3}
\end{equation}
であり、これらの単位では \(16\) と \(17\) のあいだにあり、数値としては \(16.32\,GM/c^3\) である。二つの空洞の帰還はその両側に、約 \(4.9\,GM/c^3\) と約 \(4.6\,GM/c^3\) だけ落ちる。

凍結の内側では、被積分関数は定数 \(1/A_{\min}\) である。近点が中心である動径光線は、\(r_\star\) から中心へ往復するあいだに
\begin{equation}
\label{eq:interior-radial}
\Delta t_{\mathrm{rad}}=\frac{2r_\star}{c\,A_{\min}}
  =\frac{4}{A_{\min}(1-A_{\min})}\,\frac{GM}{c^3}
\end{equation}
を費やす。一軸の飽和では \(A_{\min}<1/20\) であるから、これはすでに \(80\,GM/c^3\) を超え、三軸の上界では \(A_{\min}<1/100\) であるから \(400\,GM/c^3\) を超える。光子球をわずかに越えただけの光線は、より早く、\(b_c\sqrt{A_{\min}}\) へ下から近づく半径で転向し、その内部の往復はより短い。それでも光子の半周より長い。一軸の飽和では \(20\,GM/c^3\) より大きく、三軸の上界では \(280\,GM/c^3\) より大きい。これらの床は鋭くない。\(b_c\) にどれほど近くても、戻るあらゆる光線を光子の時計の向こう側へ置くには、十分に大きい。古典チャートには、このような有限な内部がない。

第\ref{sec:chart}節の節は、この有限な区間の上に座る。その無限個が中心へ向かって集積し、区間全体を横切る座標時間は有限のままである。凍結した速さ \(dr/dt=c A_{\min}\) そのものは零へ向かわないからである。節は帰還の上の障壁であり、第二の地平ではない。

\section{弱いレンズとしての節}
\label{sec:shells}

\(\gamma_s\) の平坦部は指数関数的に静かであるから、それらが寄与する滑らかな柱は無視する。各節は障壁である。本節の模型は、錐が許す有限なスパイクを、面積半径 \(R\)、質量 \(m\) の薄い球殻へ潰す。質量 \(m\) は負でありうる。引力の障壁であるか、斥力の障壁であるかである。錐は \(m\) を固定しない。

インパクトパラメータ \(b<R\) の視線は殻に二度出会う。射影柱は、理想的な薄い極限ではリムにおいて \(2\sigma R/\sqrt{R^2-b^2}\) のように発散する。円筒の内部の質量は発散しない。\(b\le R\) では
\begin{equation}
\label{eq:mcyl}
M_{\mathrm{cyl}}(b)=\frac{m}{R}\Bigl(R-\sqrt{R^2-b^2}\Bigr)
\end{equation}
であり、\(b\ge R\) では \(M_{\mathrm{cyl}}(b)=m\) である。弱場の偏向は、相対論の因子 \(4G/c^2\) をつけて \(\alpha(b)=4G\,M_{\mathrm{cyl}}(b)/(c^2 b)\) である。殻の外ではこれは点質量の法則である。軸上では、球対称が要求するとおり \(\alpha\to 0\) となる。\(M_{\mathrm{cyl}}\) の傾きは \(b\to R^-\) で発散する。

殻の上に座る質量のアインシュタイン半径を \(b_m\) とし、\(c=(b_m/R)^2\) と置く。その質量が負ならば \(c\) も負である。殻の角半径を単位にして、殻の内部の接線比と動径の傾きは
\begin{equation}
\label{eq:shell-profiles}
\frac{\bar\alpha}{\theta}=\frac{c}{1+\sqrt{1-\xi^2}},
\qquad
\frac{d\bar\alpha}{d\theta}
  =\frac{c}{\sqrt{1-\xi^2}\,\bigl(1+\sqrt{1-\xi^2}\bigr)}
\end{equation}
である。軸上ではどちらも \(c/2\) に等しい。外へ向かうにつれ、接線比は二倍になるだけであり、リムで有限値 \(c\) に達する。動径の傾きは有限にとどまらない。引力の殻では \(+\infty\) へ、斥力の殻では \(-\infty\) へ走る。リムで発散するのは動径固有値である。

引力の殻では、二つの比較は \(b_m\) で測られる窓に落ちる。内部の動径臨界曲線がただ一つあるのは \(R>b_m/\sqrt{2}\) のときちょうどであり、内部の接線臨界曲線がただ一つあるのは
\begin{equation}
\label{eq:shell-window}
\frac{b_m}{\sqrt{2}}<R<b_m
\end{equation}
のときちょうどである。その窓では接線曲線は動径曲線の外側にある。殻の外では、写像はその殻自身の質量の点質量である。\(R<b_m\) ならばアインシュタイン環があり、動径臨界曲線は決してない。斥力の殻は接線比も内部の傾きも負に保つので、その円盤はどちらの曲線も含まない。

殻の上の質量は未定であるから、窓は像の個数ではなく、判定条件である。節の比~\eqref{eq:radius-ratio} は節の比のままである。隣り合う臨界曲線がそれに立つ義務はない。リムを錐が要求する有限の幅へ平滑化すると、無限の動径傾きは急峻なピークで置き換わる。接線曲線はリムの上へ移らない。その曲線は、囲まれた質量の比が \(1\) に等しい場所にとどまる。

積層が変えるのは、横切られている殻の動径発散を消すことではなく、有限な内部質量の計上である。\(\gamma_s\) の符号は節から節へと交替する。臨界曲線の種類がそれに伴って交替する必要はない。光子環の列は、まったく第三の配置である。単一の曲線 \(b_c\) の相次ぐ半周であり、半周ごとに \(\mathrm{e}^{-\pi}\) で減光され、その一つの曲線へ収束する~\cite{Johnson2020}。\(b_c\) への指数的な集積と、比~\eqref{eq:radius-ratio} に入る節の対と、~\eqref{eq:shell-window} によって固定される薄い殻の臨界曲線の対とは、天空上の光の三つの配置である。本稿は、\(r_{\mathrm{ph}}\) で巻き、かつ障壁を横切る光線を積分しない。

\section{すでに見えているもの}
\label{sec:observations}

太陽の試験は外縁類に座す。光球をかすめる光線のインパクトパラメータは、\(r_s\) より \(10^5\) の程度の因子だけ大きく、原子核へ入らない。偏向は太陽の縁における弱場の法則であり、遅延は同じ因子 \(A\) のシャピロ遅延である。カッシーニの追跡は
\begin{equation}
\label{eq:cassini}
\gamma=1+(2.1\pm 2.3)\times 10^{-5}
\end{equation}
を与える~\cite{Bertotti2003}。チャートは \(\gamma=1\) の場合である。測定は、外縁の曲がりと遅延を一緒に縮尺し直す補正を制限する。光線が出会わない内部を制限はしない。

銀河の強いレンズは、より弱い場における同じ類である。アインシュタイン角と準地平の角サイズは \(\theta_E^2\sim(r_s/D_L)(D_{LS}/D_S)\) によって関係するので、その比はラジアンで測った角の平方根であり、いかなる銀河レンズについても \(1\) より何桁も小さい。\(r_\star\) にある構造は、それらの観測の精度において、アインシュタイン環も、像の位置も、時間遅延も動かさない。二つの赤方偏移にある二つの光源が形成する二重アインシュタイン環は、外縁類の環が二つである。単一クェーサーの像のあいだのフラックス比の異常も同様に外縁類の問いである。本稿はそれをねじれ節へ割り当てない。

M87* の放射環は直径 \(42\pm 3\,\mu\mathrm{as}\) をもち、フラックス比 \(\gtrsim 10{:}1\) の中心減光を囲む。その環とともに採用される角重力半径は \(\theta_g=3.8\pm 0.4\,\mu\mathrm{as}\) であり、臨界曲線~\eqref{eq:photon-sphere} の角直径は \(6\sqrt{3}\,\theta_g\) である~\cite{EHT2019}。\(6\sqrt{3}\approx 10.392\) であるから、その直径は約 \(39.5\pm 4.2\,\mu\mathrm{as}\) であり、\(42\pm 3\,\mu\mathrm{as}\) と重なる。公表された環は臨界曲線そのものではなく、臨界曲線のまわりのレンズされた放射であり、質量の推論はすでにその程度の影の縮尺を用いている。重なりは、光子球が外縁チャートの置く場所にあることの確認である。地平の検出ではない。地平は準地平が置き換えるものであり、臨界曲線はその置き換えを参照しない。

Sgr A* は、第二の質量における同じ比較である。放射環の直径は \(51.8\pm 2.3\,\mu\mathrm{as}\) である~\cite{EHT2022I}。一般相対論的シミュレーションへ較正された角重力半径は \(4.8^{+1.4}_{-0.7}\,\mu\mathrm{as}\) であり~\cite{EHT2022IV}、その区間の \(6\sqrt{3}\) 倍は約 \(43\) から \(64\,\mu\mathrm{as}\) へ走る。測定された環はその区間の内部にある。どちらの像も単一の環である。その一致が未決のまま残すものは、第\ref{sec:return}節の帰還である。

\section{予測が分かれる場所}
\label{sec:discriminants}

第\ref{sec:observations}節の一致は、接合したチャートをシュヴァルツシルト・レンズから分けない。それらを作る光線は \(r_\star\) の外にとどまるので、\(2b_c/D_A\) はシュヴァルツシルトの直径であり、カッシーニの遅延と銀河の環は、インパクトパラメータが \(r_s\) のはるか外にある。地球基線上で次世代の配列が分離できると期待するのは、単一の曲線 \(b_c\) の \(n=0\) と \(n=1\) の部分環までであり、その分離自体がすでに超分解を要する~\cite{Johnson2023}。その分離は \(4/(5\pi)\) で抑えられた比を測らない。

影の内部の明るさも、両者を分けない。貫入光線は戻るので、内部は空である必要がなく、光学的に厚いプラズマはそれを暗く保ちうる。M87* のコントラスト \(\gtrsim 10{:}1\) は和を制限する。入射パワーを熱化する面、あるいは完全に反射する面は、さらなる対象であり、Sgr A* のスペクトルと像はその二つを不利にする~\cite{EHT2022VI}。凍結 \(A=A_{\min}\) はそのいずれでもない。

接合したチャートが加えるのは、古典チャートが戻さない光線に対する、有限なキリング時間である。空洞の部分~\eqref{eq:echo-delay} はリングダウンの一周期であり、\(11.44\,GM/c^3\) または \(20.96\,GM/c^3\) であって、光子の半周~\eqref{eq:half-orbit} の傍らに置かれる。内部の部分は、それより大きい。光子球を通過するあらゆる光線は、戻り始める前に、一軸の凍結の内側で \(20\,GM/c^3\) より長く、三軸の凍結の内側で \(280\,GM/c^3\) より長く過ごす。動径光線は \(80\) より長く、\(400\) より長く過ごす。質量 \(4\times 10^6\,M_\odot\) では \(GM/c^3\approx 20\,\mathrm{s}\) であり、半周は数分、内部の床はそれより長い。第\ref{sec:observations}節の M87* の質量では、数日である。数 \(GM/c^3\) まで時刻を決めたフレアは、原理上、一つの帰還を光子の半周から分けうる。そのような分離は手元にない。

同じ間隔は、遅い重力波エコーが予測ではない理由を示す。エコー探索は、地平までの固有距離が微小な反射面を制限し、リングダウンのずっと後のパルス列を探す。ここで計算した往復は、そのリングダウンの一周期、あるいは数周期であり、その列ではない。遅い列の非検出は、\(r_\star-r_s\) ほどの幅の空洞と、速さ \(c A_{\min}\) で横切られる内部とが、期待させるものである。空洞を確認はしない。チャートは計量係数 \(A\) を供給する。重力波の境界条件は供給せず、殻の質量は未定のままである。

近点リムは、より弱い標的のままである。それらが無限遠から到達可能であるのは、面積半径が \(b_c\sqrt{A_{\min}}\) より下にあるときに限る。その切断の外の節はあらゆる貫入光線に横切られ、別々のかすめとしては現れない。その内の節は面積比~\eqref{eq:radius-ratio} でかすめられる。上界が選ばれた後に限り、かつリムが目に見えるほど明るいときに限る。

\section{結論}
\label{sec:conclusion}

ねじれ星の外縁は、シュヴァルツシルト・チャートが光をレンズするのと同じ仕方で光をレンズする。光子球は記録された両方の上界の外にあり、臨界インパクトパラメータは \(b_c=3\sqrt{3}\,GM/c^2\) のままであり、全質量のアインシュタイン環は弱場で生き延びる。太陽の遅延、銀河の環、M87* と Sgr A* の放射環が、それらの測定が実際にもつ精度において理論と会うのは、そのためである。

内部はブラックホールと会わない。\(b<b_c\) の光線は古典の外縁に転向点をもたず、\(r_s\) へ降りるキリング時間は無限であり、\(r_s\) の内側では計量係数が符号を変えている。凍結はその係数を正に保つ。同じ光線は \(r=b\sqrt{A_{\min}}\) で、\(r_\star\) の内側において転向し、その訪問の最短のものさえ、キリング時間は光子の半周を超える。節は、その有限な帰還の上にある。符号つきの薄い殻と読むと、それらは臨界曲線を加えうる。\(R>b_m/\sqrt{2}\) ならば内部の動径曲線を、窓 \(b_m/\sqrt{2}<R<b_m\) の内側に限って内部の接線曲線を加える。斥力の殻は、自身の円盤の上にどちらの曲線も加えない。各殻の質量は定まらないので、模型はどの窓が占有されるかも、倍率も、光子球で巻きかつ障壁を横切る光線の偏向も、予測しない。

したがってチャート自身が加える間隔は、第二の環ではなく、時間である。一般相対論はその光線を戻さない。ねじれ星は、凍結の内側からそれを戻す。その時計は光子の半周ではなく、遅いエコー列でもない。

\begin{thebibliography}{99}
\raggedright

\bibitem{dst-main}
Anonymous,
「双対時空の間のねじれとしての重力：一般相対性理論の双四元数的再定式化」 (2025).
\url{https://github.com/hypernumbernet/dual-spacetime-doc/blob/main/double-spacetime-theory.pdf}.

\bibitem{NarayanBartelmann1996}
R.~Narayan and M.~Bartelmann,
``Lectures on Gravitational Lensing,''
arXiv:astro-ph/9606001 (1996).

\bibitem{Bertotti2003}
B.~Bertotti, L.~Iess, and P.~Tortora,
``A test of general relativity using radio links with the Cassini spacecraft,''
\textit{Nature} \textbf{425}, 374--376 (2003).

\bibitem{EHT2019}
Event Horizon Telescope Collaboration,
``First M87 Event Horizon Telescope Results. I. The Shadow of the Supermassive Black Hole,''
\textit{Astrophys.\ J.\ Lett.} \textbf{875}, L1 (2019).

\bibitem{EHT2022I}
Event Horizon Telescope Collaboration,
``First Sagittarius A* Event Horizon Telescope Results. I. The Shadow of the Supermassive Black Hole in the Center of the Milky Way,''
\textit{Astrophys.\ J.\ Lett.} \textbf{930}, L12 (2022).

\bibitem{EHT2022IV}
Event Horizon Telescope Collaboration,
``First Sagittarius A* Event Horizon Telescope Results. IV. Variability, Morphology, and Black Hole Mass,''
\textit{Astrophys.\ J.\ Lett.} \textbf{930}, L15 (2022).

\bibitem{EHT2022VI}
Event Horizon Telescope Collaboration,
``First Sagittarius A* Event Horizon Telescope Results. VI. Testing the Black Hole Metric,''
\textit{Astrophys.\ J.\ Lett.} \textbf{930}, L17 (2022).

\bibitem{Johnson2020}
M.~D.~Johnson et al.,
``Universal interferometric signatures of a black hole's photon ring,''
\textit{Sci.\ Adv.} \textbf{6}, eaaz1310 (2020).

\bibitem{Johnson2023}
M.~D.~Johnson et al.,
``Key Science Goals for the Next-Generation Event Horizon Telescope,''
\textit{Galaxies} \textbf{11}, 61 (2023).

\end{thebibliography}

\end{document}
